In vehicular crowdsensing, a location-privacy mechanism may need to establish compliance with a spatial policy without revealing a vehicle's coordinates. Yet positioning errors, environmental effects, and synchronization uncertainty often leave only interval-valued spacetime observations. Even authenticated sources may share systematic bias or behave in a Byzantine manner. We propose a privacy-preserving compliance certification method based on signed spacetime uncertainty evidence. Given an upper bound $f$ on Byzantine issuers and an allowance $r$ for honest issuers whose intervals miss the true event, the method constructs the feasible set of states supported by at least $k=n-f-r$ issuers. A report is accepted only when this set is nonempty and wholly contained in the authorized region. The resulting decision distinguishes valid, invalid, ambiguous, and inconsistent evidence. We implement two Groth16/BN254 circuits that verify EdDSA/MiMC signatures and bind the proof to the task, issuer roster, report, threshold, challenge nonce, and expiry. Gateway tests reject tampering, context substitution, expired challenges, and cross-challenge replay; only one of eight concurrent copies of a valid request is accepted within a single gateway instance. Across 12 synthetic conditions satisfying the configured fault and coverage assumptions, mean decisive coverage is 87.52\%, compared with 75.63\% for a fixed-margin baseline. Tested subset circuits produce 164-byte proofs with median proving times of 163--255 ms on an Apple M1. The implementation is an evidence-certification layer, not a complete vehicular crowdsensing platform; its physical interpretation still depends on trusted issuer enrollment and effective sensor calibration.